\section{Related work}
\label{sec:related_work}

\subsection{Financial dependencies and econophysics}

Econophysics studies financial markets using concepts and methods from statistical physics, nonlinear dynamics, complex systems and network science. A central observation is that financial returns exhibit stylized facts that are difficult to reconcile with independent Gaussian models. Heavy tails, volatility clustering, slow decay in absolute-return autocorrelations, intermittent bursts of activity and collective co-movement patterns are recurrent findings across asset classes, countries and sampling frequencies \cite{mantegna_stanley_1999_book,cont_2001_empirical}. These regularities motivate methods that focus not only on the marginal distribution of returns, but also on dependency, synchronization and network structure.

Financial dependency analysis lies at the intersection of econophysics and financial economics. In financial economics, dependence is central to portfolio risk, diversification, contagion, factor structure and systemic vulnerability. In econophysics, the same dependence is often represented through correlation matrices, eigenspectra, filtered graphs and hierarchical structures. Raddant and Di~Matteo \cite{raddant_dimatteo_2023} emphasize that correlations, causality measures, volatility spillovers and networks offer additional views of financial systems. We follow this perspective: the goal is not merely to estimate pairwise correlations, but to decompose the correlation matrix, filter its noisy components, represent the remaining structure as networks and evaluate whether the extracted network structure has forecasting relevance.

Cross-correlation dynamics during crises have received attention because market stress often increases synchronization across assets. Zheng et~al. \cite{zheng_2012_changes} showed that changes in cross-correlation structure can serve as early indicators of systemic instability. Billio et~al. \cite{billio_2012_econometric} proposed econometric measures of connectedness and systemic risk in financial sectors. Junior and Franca \cite{junior_2012_correlation} documented that correlations among global stock indices tend to increase during crises, reducing the effective dimensionality of international markets. These findings motivate the rolling-correlation and crisis-synchronization analysis undertaken in this paper, where market-wide co-movement is treated as an object of study instead of a nuisance.

The relevance of these studies for the present manuscript is twofold. First, they show that financial dependence is dynamic and crisis-sensitive. Second, they indicate that a single unfiltered correlation matrix may conceal different layers of structure: a broad market component, group or sectoral organization and residual noise. This motivates the use of Random Matrix Theory and filtered networks as additional tools for separating and interpreting these layers.

\subsection{Random Matrix Theory in financial correlations}

Random Matrix Theory entered finance as a way to separate information from noise in large correlation matrices. Laloux et~al. \cite{laloux_1999_noise,laloux_2000_random} showed that a substantial part of the eigenspectrum of financial correlation matrices can resemble the spectrum predicted for random matrices. The relevant benchmark is the Mar\v{c}enko--Pastur distribution \cite{marcenko_1967_distribution}, which sets bounds for the eigenvalues of a random correlation matrix estimated from independent standardized time series. Plerou et~al. \cite{plerou_2002_random} further studied eigenvalues and eigenvectors outside the random band and read the largest eigenvalue as a market-wide mode.

The intuition is that the empirical correlation matrix contains both structured and noisy components. If $N$ assets are observed over $T$ periods, the ratio $Q=T/N$ determines the width of the random-matrix noise band. Eigenvalues inside this band are treated as compatible with sampling noise, whereas eigenvalues above the upper edge are candidates for collective factors. This interpretation is useful in financial markets because broad market co-movement can dominate the first eigenvalue, while weaker sectoral or group modes may appear in subsequent eigenvectors.

Subsequent work extended RMT-based methods for covariance and correlation cleaning. Bun et~al. \cite{bun_2017_cleaning} reviewed and unified several approaches for cleaning large covariance matrices using random matrix theory, connecting theoretical results with practical financial applications. Guhr et~al. \cite{guhr_muller_1998} give a broader mathematical review of random matrix theory and its applications in complex systems. In finance, RMT has been used not only for noise filtering, but also for portfolio optimization, factor identification and the study of collective market behaviour.

Our RMT interpretation is cautious. The largest eigenvalue captures a market mode because it reflects broad positive co-movement across many assets. Eigenvalues above the Mar\v{c}enko--Pastur upper edge, excluding the first one, are read as group or sector candidates. Eigenvalues inside the random band are treated as noisy residual components. These labels should be read as filtering categories; they do not imply that each eigenvector has a unique economic interpretation. We use this spectral separation as the input for network construction and volatility forecasting in the Brazilian equity market.

\subsection{Financial networks: MST, PMFG and filtered graphs}

Financial networks transform a dependency matrix into a graph representation. The usual approach begins by converting correlations into distances, so that strongly correlated assets are close in the induced metric space. Mantegna \cite{mantegna_1999_hierarchical} introduced the Minimum Spanning Tree (MST) as a way to reveal hierarchical organization in financial markets. The MST retains the connected tree with minimum total distance and therefore contains exactly $N-1$ edges for $N$ assets. Its main advantage is interpretability: it retains a sparse backbone of the strongest dependencies. Its main limitation is that it removes cycles and many potentially meaningful local connections.

Subsequent studies applied MST-based methods to study market taxonomy, instability and dynamics. Onnela et~al. \cite{onnela_2003_dynamics} analysed the evolution of MST-based market correlations and their implications for portfolio structure. Bonanno et~al. \cite{bonanno_2004_networks} extended the network approach to multiple world markets, showing that MSTs often organize assets according to sectoral and geographical structure. Song et~al. \cite{song_2012_evolution} studied the evolution of worldwide stock markets through correlation-based graphs, documenting time variation in network structure.

The Planar Maximally Filtered Graph (PMFG) relaxes the sparsity of the MST while preserving planarity \cite{tumminello_2005_tool}. For $N$ nodes, a PMFG contains $3N-6$ edges and includes the MST as a subgraph when both are constructed from the same ranked distance matrix. Because the PMFG can retain triangles and 4-cliques, it is more suitable for studying local connectivity, clustering, hub reorganization and group-level structure. Tumminello et~al. \cite{tumminello_2010_correlation} further explored the relationship between correlation, hierarchies and networks in financial markets. Massara et~al. \cite{massara_2016_tmfg} later introduced the Triangulated Maximally Filtered Graph (TMFG), which extends this class of planar filtering methods and improves scalability for large systems.

Filtered networks are relevant for risk analysis. Pozzi et~al. \cite{pozzi_2013_spread} used PMFG-based analysis to show that peripheral assets in financial networks can offer better diversification properties than central hubs. Aste et~al. \cite{aste_2010_correlation} studied correlation dynamics in volatile markets using filtered network approaches. Kenett et~al. \cite{kenett_2012_dominating} used partial correlation analysis to reveal the dominant role of the financial sector in stock market networks. These studies motivate the hub-reorganization analysis presented here. By comparing original and RMT-filtered networks, we ask whether centrality is driven mainly by market-wide co-movement or whether localized central nodes emerge after removing the dominant market mode.

\subsection{Financial networks in emerging markets}

Network methods have also been applied to emerging markets, where market structure can differ substantially from that of large developed equity markets. Emerging markets may combine concentrated sectoral composition, macroeconomic sensitivity, lower liquidity in parts of the market and exposure to global risk cycles. These features make them attractive systems for dependency analysis because the correlation structure may reflect both domestic and international channels.

Tabak et~al. \cite{tabak_2010_topology} studied the structure of the Brazilian interbank network, documenting structural properties and vulnerability to targeted failures. Silva et~al. \cite{silva_2016_structure} analysed the structure and dynamics of the global financial network, with attention to how emerging economies connect to the broader system. Wang et~al. \cite{wang_2018_correlation} studied the correlation structure of world stock markets using Pearson and partial-correlation networks, showing that emerging markets can display distinctive dependency patterns.

Studies focused more directly on Brazilian markets include Pereira et~al. \cite{pereira_2019_multiscale}, who applied detrended cross-correlation analysis to construct multiscale networks of stock markets including B3, and Leonel Caetano and Yoneyama \cite{leonel_2022_brazilian}, who applied complex network analysis to the Brazilian stock market. Stosic et~al. \cite{stosic_2015_correlations} studied correlations and cross-correlations in Brazilian agrarian commodities and stocks. These studies show that Brazilian financial data contain rich structural properties, but the literature remains fragmented across interbank networks, commodities, cross-market relations and equity-market structure.

The present paper contributes to this literature by focusing on a broad B3 equity universe and by connecting spectral filtering, planar networks and volatility forecasting in a single design. This integration matters because a network's interpretation depends on the matrix used to build it. A network constructed from the original correlation matrix may be dominated by the market mode, while a network constructed from group-mode components may reveal weaker but more localized dependency channels. This distinction is especially relevant in emerging markets, where common macroeconomic exposure and sectoral heterogeneity coexist.

\subsection{Volatility forecasting and machine learning}

Volatility forecasting has increasingly incorporated machine-learning methods because realized-volatility targets can be combined with large feature sets, nonlinear interactions and strict temporal validation. The theoretical foundations of realized-volatility measurement were established by Andersen et~al. \cite{andersen_2003_modeling}, and Poon and Granger \cite{poon_granger_2003} give an influential review of volatility forecasting methods. In the present study, the volatility-modelling component is retrospective: it asks whether features extracted ex post from RMT and filtered networks are associated with supervised-model performance for future realized volatility. It does not claim a point-in-time forecasting advantage from these full-sample structural descriptors.

Forecast evaluation is delicate because realized volatility is an observable proxy for latent conditional variance. Loss functions can rank models differently depending on their robustness to measurement error and their treatment of volatility underestimation. QLIKE is widely used because of its robustness properties under noisy volatility proxies \cite{patton_2011_volatility}. For this reason, the forecasting experiment in this paper evaluates models using multiple metrics, including MAE, RMSE, QLIKE and out-of-sample $R^2$.

Machine-learning methods offer flexible nonlinear approximations and can incorporate large feature sets. Random Forests \cite{breiman_2001_random}, gradient boosting \cite{friedman_2001_greedy} and regularized linear models \cite{hoerl_kennard_1970_ridge} can combine lagged volatility, rolling statistics, market-wide variables and network centrality measures. Bucci \cite{bucci_2020_realized} applied neural networks to realized-volatility forecasting, showing that flexible models can capture nonlinear volatility dynamics. Christensen et~al. \cite{christensen_2022_realized} developed a machine-learning framework for volatility forecasting and emphasized the importance of careful feature engineering. Zhang and Broadstock \cite{zhang_2023_network_volatility} explored how financial network structure relates to market risk from a machine-learning perspective.

In financial forecasting, however, flexibility can amplify noise and produce unstable gains if temporal validation is not strict. For this reason, we use a chronological train-validation-test split and treat network features as potential additional predictors, not replacements for lagged volatility information. The nested feature-set design follows an ablation logic: Set~A measures the predictive value of standard volatility information, Set~B tests the added value of market and RMT variables, and Set~C tests whether MST and PMFG structure adds further information beyond these baseline predictors. Simple MLP and CNN-1D models are included as incremental neural extensions rather than as an exhaustive architecture search.
